\section{Conformance performance and residual uncertainty}\label{sec:C19:evaluation}
Protocol evaluation should separate schema conformance from application correctness. Conformance tests can validate required metadata, unsupported versions, capability errors, pagination, and response correlation. They can inspect whether the current modern request model is implemented rather than a historical initialization flow. They cannot establish that an agent's chosen action is useful or that a returned research artifact is scientifically sound.

Use a version matrix with modern-only, legacy-only, and dual-era MCP peers. Include missing required metadata, declared but unsupported capabilities, an unknown version, and a recognized modern rejection. Test whether fallback follows the binding's rules instead of treating every error as a reason to initialize a legacy session. Preserve exact specification and SDK revisions in the result. A package release number alone is insufficient when examples or generated schemas change independently.

For A2A, test direct messages, stateful tasks, interrupted work, terminal outcomes, and artifact updates. Include delayed and duplicated updates, failed streams, and concurrent follow-ups. Check that the application preserves task ownership and does not interpret a progress event as completion. Independently verify a sample of final artifacts against the requested predicate. Protocol-level terminal status and task-level validity should remain separate metrics.

Performance accounting includes discovery, schema transfer, serialization, credential flows, network latency, polling or streaming overhead, model decisions, and remote work. Compare warm and cold operation, and identify which metadata is cacheable under which principal. A cached catalog can reduce overhead while becoming stale after a permission or tool change. Latency gains should therefore be reported alongside freshness and correctness behavior.

The official MCP architecture documentation identifies an additional current-version change: sampling and some prior utilities are deprecated in the 2026-07-28 era, and change notifications are delivered through opt-in subscriptions~\cite{mcp2026architecture}. Implementations and surveys should not infer present support from an older tutorial. Deprecation also differs from immediate removal; the exact feature lifecycle must be checked before asserting that a capability is unavailable.

No runtime conformance or cross-protocol benchmark has been executed for this chapter. The normative facts come from official specifications inspected at the cutoff; the execution trace, failure analysis, and test design are explicitly deductions and proposals. The main technical conclusion is bounded: interoperable messages make component composition possible, while durable state, authority, evidence quality, and verified outcomes remain contracts of the complete agent system. Those contracts determine whether interoperability improves reliable autonomy.
